\newcommand{\PRED}{\ensuremath{\mbox{pred}}}
\newcommand{\FACT}{\ensuremath{\mbox{fact}}}
\newcommand{\ID}{\ensuremath{\mbox{I}}}
\newcommand{\IF}{\ensuremath{\mbox{if}}}
\newcommand{\ADD}{\ensuremath{\mbox{add}}}
\newcommand{\MULT}{\ensuremath{\mbox{mult}}}
\newcommand{\ISZERO}{\ensuremath{\mbox{iszero}}}
\newcommand{\CN}[1]{\ensuremath{\underline{#1}}}
\newcommand{\Y}{\ensuremath{Y}}
\newcommand{\ZERO}{\CN{0}}
\newcommand{\ONE}{\CN{1}}
\newcommand{\TWO}{\CN{2}}
\newcommand{\LAM}[2]{\ensuremath{\lambda #1.#2}}


\section{Homework 10: Untyped Lambda Calculus}

This homework contains two exercises on the untyped lambda calculus. For readability,
we will use the following definitions from the lecture:

\begin{align*}
\ID &= \LAM{x}{x} &\top &= \LAM{xy}{x} &\bot &= \LAM{xy}{y} \\
\IF &= \LAM{bxy}bxy &\ZERO &= \LAM{Sx}{x} &\ONE &= \LAM{Sx}{Sx} \\
\TWO &= \LAM{Sx}{S(Sx)} &\underline{n} &= \LAM{Sx}{S^n(x)} &\ADD &= \LAM{mn}{mSn} \\
\MULT &= \LAM{mn}{m(\ADD{}\, n)\,\ZERO}  &\ISZERO&=\LAM{n}{n(\LAM{x}{\bot})\top} &\Y &= \LAM{f}{(\LAM{x}{fxx})(\LAM{x}{fxx})}
\end{align*}

\subsection{Predecessor function for Church numerals}

\[
\PRED = \LAM{nfx}{n(\LAM{gh}{h(gf)})(\LAM{u}{x})(\LAM{u}{u})}
\]

% Below is a derivation and explanation of the predecessor (\(\PRED\)) function for Church numerals, followed by the final definition in a copyable LaTeX code block. We begin by recalling what the Church numeral \(\underline{n}\) and its successor function \(\SUCC\) look like, and then proceed to construct a function that, when applied to a numeral \(n\), returns \(n-1\).

The predecessor function \(\PRED\) should take a numeral \(n\) and return \(n-1\). We know that:
\begin{itemize}
  \item \(\PRED(\ZERO) = \ZERO\), since the predecessor of zero should not go below zero.
  \item For any \(n > 0\), \(\PRED(n) = n - 1\).
\end{itemize}

A clever technique to define \(\PRED\) uses a ``pairing'' function that keeps
track of two values as we iterate through the numeral. The idea is to transform
the iteration function of the numeral into one that carries along a state
representing the current and previous "count". By the time we finish applying
the numeral's function the appropriate number of times, we end up with the
predecessor.

We want an iteration that, given a function \(f\) and an initial value \(x\),
produces something related to \((n-1)\). Consider we apply \(n\) times a
function that takes a pair \((g, h)\) and returns a new pair that effectively
shifts \(f\) from ``next'' to ``previous'':

\begin{enumerate}
  \item Start with a pair:
  \begin{itemize}
    \item A ``dummy'' function \(\LAM{u}{x}\) that, if applied, yields \(x\).
    \item The identity function \(\ID = \LAM{u}{u}\) which essentially acts as the successor chain.
  \end{itemize}
   We think of this initial pair as \((\LAM{u}{x}, \LAM{u}{u})\). The idea is
   that the second element of the pair eventually ends up ``one step ahead'' of
   the first. After iterating \(n\) times, the first component of the pair will
   encode the predecessor computation.
  \item Each iteration step transforms \((g,h)\) into \((h, \LAM{z}{h(gf)})\) or
something equivalent. More concretely, we define a function \(\LAM{gh}{h(gf)}\)
that, given a pair \((g,h)\), produces a new function related to this shifting.
By repeatedly applying this, we ``shift'' through the functions, so that after
\(n\) steps, the ``front'' shifts out one step ahead, and what remains in the
first component is effectively the predecessor function applied to \(f\) and
\(x\).
\end{enumerate}


This yields the following lambda term for \(\PRED\):

\[
\PRED = \LAM{nfx}{n(\LAM{gh}{h(gf)})(\LAM{u}{x})(\LAM{u}{u})}.
\]

We illustrate the correctness of this definition by $\beta$-reducing \(\PRED\,\ZERO\) and \(\PRED\,\TWO\):

\subsubsection*{Reduction of \(\PRED\,\ZERO\)}

\[
\begin{aligned}
\PRED\,\ZERO &= (\LAM{n f x}{n(\LAM{g h}{h(g f)})(\LAM{u}{x})(\LAM{u}{u})})\,\LAM{Sx'}{x'} \\[4pt]
&\rightarrow_\beta \LAM{f x}{(\LAM{S x'}{x'})(\LAM{g h}{h(g f)})(\LAM{u}{x})(\LAM{u}{u})} \\[4pt]
&\rightarrow_\beta \LAM{f x}{(\LAM{x'}{x'})(\LAM{u}{x})(\LAM{u}{u})} \\[4pt]
&\rightarrow_\beta \LAM{f x}{(\LAM{u}{x})(\LAM{u}{u})} \\[4pt]
&\rightarrow_\beta \LAM{f x}{x} \\
&= \ZERO{}
\end{aligned}
\]

\subsubsection*{Reduction of \(\PRED\,\TWO\)}

\[
\begin{aligned}
  \PRED\,\TWO
  &= (\LAM{n f x}{n(\LAM{g h}{h(g f)})(\LAM{u}{x})(\LAM{u}{u})})\,(\LAM{S x}{S(S x)}) \\[6pt]
  &\rightarrow_\beta \LAM{f x}{(\LAM{S x'}{S(S x')})(\LAM{g h}{h(g f)})(\LAM{u}{x})(\LAM{u}{u})} \\[6pt]
  &\rightarrow_\beta \LAM{f x}{(\LAM{x'}{(\LAM{g h}{h(g f)})((\LAM{g h}{h(g f)}) x')})(\LAM{u}{x})(\LAM{u}{u})} \\[6pt]
  &\rightarrow_\beta \LAM{f x}{(\LAM{g h}{h(g f)})((\LAM{g h}{h(g f)})(\LAM{u}{x}))(\LAM{u}{u})} \\[6pt]
  &\rightarrow_\beta \LAM{f x}{(\LAM{g h}{h(g f)})(\LAM{h}{h(x)})(\LAM{u}{u})} \\[6pt]
  &\rightarrow_\beta \LAM{f x}{(\LAM{h}{h((\LAM{h}{h(x)}) f)})(\LAM{u}{u})} \\[6pt]
  &\rightarrow_\beta \LAM{f x}{(\LAM{h}{h(f(x))})(\LAM{u}{u})} \\[6pt]
  &\rightarrow_\beta \LAM{f x}{f(x)} \\[6pt]
  &= \ONE
\end{aligned}
\]

\subsection{$\beta$-Reduction of $\FACT{}$ $2$}

For the sake of readability, we assume that all arithmetic functions defined in
the lecture correspond to their intended behavior (i.e., we assume that $\ADD~n
m = n + m$, $\MULT~n m = n \cdot m$, and $\ISZERO~n = \top$ if $n = 0$ and
$\bot$ otherwise). Steps following from these definitions are written using
a double arrow $\Rightarrow_\beta$.

\begin{align*}
\FACT\,\TWO
&= (\Y(\LAM{f n}{\IF(\ISZERO\, n)\ONE(\MULT\, n(f(\PRED\, n)))}))\,\TWO \\[6pt]
&= ((\LAM{f}{f(\LAM{x}{fxx})(\LAM{x}{fxx})})(\LAM{f n}{\IF(\ISZERO\, n)\ONE(\MULT\, n(f(\PRED\, n)))}))\,\TWO \\[6pt]
&\rightarrow_\beta (\LAM{x}{(\LAM{f n}{\IF(\ISZERO\, n)\ONE(\MULT\, n(f(\PRED\, n)))}) (x x)})\\&\quad~ (\LAM{x}{(\LAM{f n}{\IF(\ISZERO\, n)\ONE(\MULT\, n(f(\PRED\, n)))}) (x x)})\,\TWO \\[6pt]
&\rightarrow_\beta (\LAM{n}{\IF(\ISZERO\, n)\ONE(\MULT\, n((Y(\LAM{f n}{\IF(\ISZERO\, n)\ONE(\MULT\, n(f(\PRED\, n)))}))(\PRED\, n)))})\,\TWO \\[6pt]
&\rightarrow_\beta \IF(\ISZERO\,\TWO)\ONE(\MULT\,\TWO(\FACT(\PRED\,\TWO))) \\[6pt]
&\rightarrow_\beta \IF(\bot)\ONE(\MULT\,\TWO(\FACT\,\ONE)) \\[6pt]
&\rightarrow_\beta \MULT\,\TWO(\FACT\,\ONE)
\end{align*}

The first two reductions of computing $\FACT~\ONE$ and $\FACT~\ZERO$ are
analogous to the previous computation and will thus be omitted.

\begin{align*}
\FACT\,\ONE
&= (Y(\LAM{f n}{\IF(\ISZERO\, n)\ONE(\MULT\, n(f(\PRED\, n)))}))\,\ONE \\[6pt]
&\Rightarrow_\beta \IF(\ISZERO\,\ONE)\ONE(\MULT\,\ONE(\FACT(\PRED\,\ONE))) \\[6pt]
&\rightarrow_\beta \IF(\bot)\ONE(\MULT\,\ONE(\FACT\,\ZERO)) \\[6pt]
&\rightarrow_\beta (\bot)\ONE(\MULT\,\ONE(\FACT\,\ZERO)) \\[6pt]
&\rightarrow_\beta \MULT\,\ONE(\FACT\,\ZERO)
\end{align*}

\begin{align*}
\FACT\,\ZERO
&= (Y(\LAM{f n}{\IF(\ISZERO\, n)\ONE(\MULT\, n(f(\PRED\, n)))}))\,\ZERO \\[6pt]
&\Rightarrow_\beta \IF(\ISZERO\,\ZERO)\ONE(\MULT\,\ZERO(\FACT(\PRED\,\ZERO))) \\[6pt]
&\rightarrow_\beta \IF(\top)\ONE(\MULT\,\ZERO(\FACT\,\ZERO)) \\[6pt]
&\rightarrow_\beta (\top)\ONE(\MULT\,\ZERO(\FACT\,\ZERO)) \\[6pt]
&\rightarrow_\beta \ONE
\end{align*}

From these computations, we can conclude that $\FACT\,\TWO = \MULT\,\TWO(\FACT\,\ONE) = \MULT\,\TWO(\MULT\,\ONE(\FACT\,\ZERO)) = \MULT\,\TWO(\MULT\,\ONE\,\ONE) = \MULT\,\TWO\,\ONE = \TWO$.
